\documentclass{article}

% arXiv-standard packages
\usepackage[T1]{fontenc}
\usepackage[utf8]{inputenc}
\usepackage[margin=1in]{geometry}
\usepackage{amsmath,amssymb}
\usepackage{booktabs}
\usepackage{array}
\usepackage{hyperref}
\usepackage{url}
\usepackage{microtype}
\usepackage{parskip}
\usepackage{xcolor}

\hypersetup{
  colorlinks=true,
  linkcolor=blue!70!black,
  citecolor=blue!70!black,
  urlcolor=blue!70!black,
}

\title{\textbf{We Have Been Building Haystacks When We Need\\
Forests of Knowledge Trees}\\[0.5em]
\large Knowledge Graph Compilation as an Alternative to Universal\\Parameter-Space Compression}

\author{Eric G.\ Suchanek, Ph.D.\\
\small Independent Researcher, Flux-Frontiers\\
\small \href{mailto:suchanek@flux-frontiers.com}{suchanek@flux-frontiers.com}}

\date{May 2026}

\begin{document}

\maketitle

\begin{abstract}
Modern AI rests on a seductive premise: compress all human knowledge into a
single high-dimensional parameter space and recover it through learned
similarity.  We argue this premise is structurally flawed.  Universal
compression yields not a knowledge base but a \emph{haystack}---every fact
occupying the same undifferentiated space, entangled across incompatible
ontologies, grammars, and languages.  Retrieval from such a structure is
probabilistic archaeology, not lookup.  We present KGRAG (Knowledge Graph
Retrieval-Augmented Generation) as a fundamental architectural inversion:
rather than compressing knowledge into weights, we \emph{compile} it into
structured, domain-specific knowledge graphs organized as a federated forest of
knowledge trees.  The model's role shrinks to reader and synthesizer of
explicitly structured facts.  For domains with strong formal ontologies---Python
source code and metabolic biochemistry being the clearest examples---no model
training is required at all.  The corpus \emph{is} the knowledge.  What remains
is simply the need for a better reader.
\end{abstract}

% ============================================================
\section{The Haystack Problem}
% ============================================================

The history of modern AI is, in large part, a history of making haystacks
bigger.  A language model is trained on a corpus---books, articles, code,
scientific papers, legal documents---and billions of parameters are adjusted
until the model produces probable next tokens.  The result is a single unified
parameter space in which the biochemistry of ATP synthesis coexists with the
grammar of 14th-century Middle English poetry, Python call conventions, and the
procedural rules of maritime law.  This is not a knowledge base.  It is a
haystack.

\subsection{Ontological Collapse}

Every domain of human knowledge has its own ontology: a system of entities,
relationships, and rules that give its claims their meaning.  Biochemistry has
metabolites, reactions, enzymes, and stoichiometry.  Programming languages have
functions, call sites, modules, and type systems.  Legal codes have statutes,
provisions, cross-references, and hierarchical authority.

When these ontologies are compressed into a shared parameter space, their
boundaries dissolve.  The model learns statistical co-occurrence across all of
them simultaneously.  The result is \emph{ontological collapse}: the model
knows something about everything, and each domain's knowledge is entangled with
every other domain's.  There is no clean substrate for any one ontology---only
interference patterns across all of them.

\subsection{Retrieval is Approximate by Construction}

A language model does not look up facts.  It generates probable continuations of
input sequences.  When asked a factual question, it produces the statistically
likely answer given its training distribution.  This is not retrieval---it is
interpolation from the nearest region of a very high-dimensional space.

When the relevant fact is rare, ambiguous, contradicted by other training
examples, or expressed in domain-specific terms that collide with
surface-similar terms elsewhere, the model interpolates from the surrounding
haystack.  The result is plausible but incorrect: hallucination.  Hallucination
is not a bug.  It is the structural consequence of treating a compression
artifact as a knowledge store.

\subsection{Scaling Intensifies the Problem}

The dominant response to these limitations has been scale.  Larger models,
larger corpora, more parameters.  This reasoning is correct about capacity but
wrong about structure.  A larger haystack is still a haystack.  More parameters
provide more room for ontological entanglement, not a substrate for ontological
boundaries.  The larger the model, the harder it becomes to audit or trace any
individual claim.

The Mixture of Experts (MoE) architecture \cite{shazeer2017moe} represents a
partial recognition of this problem.  By routing inputs to specialized
sub-networks, MoE recovers something like domain-specific subspaces.  The
routing mechanism learns implicitly to separate domains that should have been
separated explicitly.  This is progress---but it is progress within the haystack
paradigm.  The fundamental substrate remains a compression of knowledge into
weights.

% ============================================================
\section{The Inversion: Forests of Knowledge Trees}
% ============================================================

The solution is not a better haystack.  It is a fundamentally different
substrate.

Human knowledge is not homogeneous.  It is organized by domain, by ontology, by
the formal structures that give each field its expressive power.  A biochemical
pathway has a different structure than a legal statute, which has a different
structure than a Python module.  These differences are not incidental.  They are
the \emph{source} of the knowledge's meaning.  The right substrate preserves
these structures rather than dissolving them.  That substrate is the knowledge
graph.

\subsection{Knowledge Graphs as Compiled Knowledge}

A knowledge graph derived from a formally structured source is not an
approximation of that source.  It is a \emph{relational representation} of
it---the same information, reorganized for navigation, traversal, and retrieval.

Consider a Python codebase---and this is not a hypothetical.  PyCodeKG
(\url{https://github.com/Flux-Frontiers/pycode\_kg}) does this today.  The
abstract syntax tree (AST) of a Python file is a formal document; its call
graph, import graph, and class hierarchy are theorems derivable by deterministic
parsers without inference.  A knowledge graph built from those parsers is correct
by construction.  Every \texttt{CALLS} edge exists because a call expression was
found at a specific AST node.  Every \texttt{INHERITS} edge exists because an
inheritance declaration was parsed.  No edge is guessed; no relationship is
inferred.  PyCodeKG demonstrates the principle at production scale---the
technique is validated, not theoretical.

The same principle holds across a broad range of domains, shown in
Table~\ref{tab:domains}.

\begin{table}[h]
\centering
\caption{Representative strong-ontology domains and their knowledge graph kinds.}
\label{tab:domains}
\begin{tabular}{lll}
\toprule
\textbf{Domain} & \textbf{Formal structure} & \textbf{Graph kind} \\
\midrule
Python source code        & Abstract syntax tree              & Code graph \\
Biochemical pathways      & Reaction schemas (KEGG, BioCyc)   & Metabolic graph \\
Protein structures        & PDB coordinate files              & Structural graph \\
Legal codes               & Statutory hierarchies             & Legal graph \\
Scientific documents      & Section/citation/entity graphs    & Document graph \\
Diary and memory          & Temporal event sequences          & Episodic graph \\
\bottomrule
\end{tabular}
\end{table}

In each case, the formal structure of the domain \emph{is} the ontology.  A
knowledge graph derived from it inherits that ontology explicitly rather than
approximating it statistically.

\subsection{The Forest Metaphor}

A single knowledge graph is a tree: a structured, navigable, rooted hierarchy of
entities and relationships.  Each domain produces its own tree.  A biochemistry
corpus produces a metabolic tree.  A codebase produces a code tree.  A legal
corpus produces a statutory tree.

These trees are not isolated.  They are related through shared entities,
cross-domain references, and questions that span multiple domains simultaneously.
A scientist studying oxidative phosphorylation must traverse the metabolic
pathway graph, the protein structure graph, the analysis code graph, and the
relevant literature graph in a single inquiry.

This is the forest: a federated collection of domain-specific knowledge trees,
each structurally sound within its own ontology, queryable together through a
single interface.  Where the haystack is undifferentiated and approximate, the
forest is structured and exact.  Where the haystack dissolves ontological
boundaries, the forest preserves them.  Where the haystack retrieves by
probability, the forest retrieves by traversal.

\subsection{The Model Becomes a Reader}

In the haystack paradigm, the model \emph{is} the knowledge.  In the forest
paradigm, the knowledge is in the graphs.  The model's role reduces to two
well-defined tasks: (1)~\textbf{query formulation}---translating a
natural-language question into a retrieval operation over the appropriate part of
the forest; and (2)~\textbf{synthesis}---reasoning over a structured,
provenance-tagged set of facts returned by graph traversal.

Both tasks are fundamentally simpler than being a knowledge store.  A model
given structured, verified facts need not memorize biochemistry, law, and Python
call conventions simultaneously.  It needs to be a good reader.  A good reader
is a far smaller, simpler, and more auditable system than a compressed
representation of all human knowledge.

% ============================================================
\section{Strong-Ontology Domains: PyCodeKG and MetaboKG}
% ============================================================

The strongest evidence for this paper's central claim comes from two existing
systems developed at Flux-Frontiers, each handling a domain with a formally
specified ontology.

\subsection{PyCodeKG: The Proof Already Exists}

Python is the ideal test case: a formal grammar, a deterministic parser, and
unambiguous semantics for every construct.  PyCodeKG exploits this fully.
Three deterministic AST passes extract every definition, call site, import, and
inheritance relationship, then data-flow edges---reads, writes, and attribute
accesses.  A post-build symbol resolution step links unresolved call stubs to
their first-party definitions.  The result is a knowledge graph with thousands
of nodes and edges, all derived from syntax, with zero model involvement in
construction.

PyCodeKG is not a prototype.  It is a production system in active use, queryable
today through KGRAG's federation layer and the Claude Code MCP interface.  It
demonstrates the core thesis at scale: \textbf{when the ontology is formal, the
graph is the knowledge, and no training is required.}

The relationship between knowledge quality and model sophistication, in this
domain, is essentially zero.  More Python source code means a larger graph and
better answers.  That is the entire scaling law.

\subsection{MetaboKG: Extending the Principle to Biochemistry}

Biochemistry has what we call a \emph{strong ontology}: a system of entities and
relationships that is not merely conventional but formally specified.  The
chemical identity of ATP is not a statistical pattern---it is a molecular
formula, a structural graph, and a set of reaction roles precisely defined by
decades of biochemical formalization.  The relationship between NADH and
Complex~I of the electron transport chain is a stoichiometric, spatial, and
mechanistic fact encoded in KEGG \cite{kanehisa2023kegg}, BioCyc, and the PDB.

MetaboKG (\url{https://github.com/Flux-Frontiers/metabo\_kg}) applies the same
compilation principle to metabolic pathway data.  A query for ``ATP synthesis in
Complex~V'' does not require approximate semantic matching---it requires
traversal from the ATP node through reaction edges to enzyme nodes to structural
evidence.  The ontology is the understanding.  The graph traversal \emph{is} the
reasoning.

\subsection{No Training Required}

This is the central claim, stated plainly: \textbf{for domains with strong
ontologies, no model training is required to answer domain-specific questions.}

A biochemical knowledge graph built from curated, formally structured databases
contains the answers to biochemical questions as paths through its structure.  A
system that can traverse those paths---even a simple rule-based traversal
engine---can answer questions that would challenge a large language model.  The
answers are correct by construction, traceable to primary sources, and immune to
hallucination.

MetaboKG's bottleneck is not model capability.  It is corpus size.  The
architecture is validated by PyCodeKG; MetaboKG applies the same pattern to a
different formal ontology.  The remaining work is expanding the graph.

\subsection{The Reader Residual}

Even in the strong-ontology case, a reader remains useful at the boundaries.
Natural-language queries must be mapped to graph entry points; ambiguous terms
must be disambiguated; synthesized answers must be rendered in prose.

But this residual reader problem is categorically different from general
knowledge compression.  Disambiguation within a well-defined ontology is a
constrained task: there is a known vocabulary of entities, and the reader's job
is to match query terms to that vocabulary.  Existing NLP tools---named entity
recognition, terminology normalization, ontology lookup---already solve most of
this without a large language model.  For synthesis, even a small model given
structured, provenance-tagged context can produce accurate, well-grounded
responses.  The task is tractable precisely because the model is given facts, not
asked to construct them.

% ============================================================
\section{The Weak-Ontology Case: NLP as a Corpus Feeder}
% ============================================================

Not all domains have strong ontologies.  Natural language text---journalism,
personal correspondence, informal technical writing---lacks the formal structure
that enables deterministic graph extraction.  For these domains the haystack
critique is weaker: there is no formal structure to exploit.

The forest paradigm nevertheless offers a path.  Dependency parsers, coreference
resolvers, named-entity recognizers, and relation extractors can impose
provisional structure on unstructured text.  The result is not as precise as a
formally structured knowledge graph, but it is substantially more structured than
parameter-space compression.  Entities are identified; relationships are typed;
temporal and spatial context is tagged.

This provisional structure, fed into the KGRAG pipeline, produces a knowledge
graph that is weaker than the strong-ontology case but far stronger than the
haystack alternative.  As the corpus grows and patterns recur, provisional
structures harden into stable graph elements.  The progression is from haystack
to forest, even in domains where the forest starts as scrubland.

% ============================================================
\section{KGRAG: Architecture of the Forest}
% ============================================================

KGRAG implements the forest paradigm as a federated system.  Its architecture
reflects the core claims of this paper.

\subsection{Domain-Specific Compilation}

Each domain is handled by a dedicated \emph{compiler}---a parser that understands
the formal structure of the domain and extracts a knowledge graph deterministically.
PyCodeKG compiles Python source code via AST analysis.  DocKG compiles Markdown
documentation via document parse trees.  MetaboKG compiles metabolic pathway data
via reaction and enzyme schemas from KEGG and BioCyc.

The compilation step is analogous to a software compiler: it transforms source
material into a lower-level representation that preserves meaning but makes it
amenable to navigation and retrieval.  The compiled artifact---a SQLite database
of nodes and edges augmented with a semantic vector index---is correct by
construction and can be rebuilt identically from the same source at any time.

\subsection{Structural Traversal with Semantic Entry Points}

KGRAG queries proceed in two phases.  First, \emph{semantic seeding}: given a
natural-language query, sentence embeddings find the $k$ nodes most semantically
related to the query terms.  Second, \emph{structural expansion}: breadth-first
traversal from those seed nodes follows typed edges through the graph, collecting
everything reachable within $h$ hops.

This separation is principled.  Semantic search handles the mapping from
natural language to graph entry points---the reader's disambiguation task.
Structural traversal handles retrieval of related knowledge---a deterministic
graph operation immune to approximation errors.  The vector index is an
accelerant; the structural graph is the authority.  This two-phase architecture
relates to, but is architecturally distinct from, GraphRAG
\cite{edge2024graphrag}, which applies graph-based community detection to LLM
outputs rather than compiling from formal sources.

\subsection{Federation: One Interface for the Whole Forest}

Every domain-specific knowledge graph exposes the same adapter interface.  The
federation layer---the KGRAG orchestrator---holds a registry of all available
graphs and dispatches every query to all of them simultaneously.  Results are
merged and globally ranked by relevance with full provenance attribution.

The scientist's query for ``oxidative phosphorylation'' fans out to the metabolic
graph, the protein structure graph, the analysis code graph, and the literature
graph simultaneously.  Every relevant result from every domain returns through a
single interface.  The forest answers as a whole.

\subsection{Provenance as a First-Class Property}

Every node in every KGRAG graph carries a stable identifier encoding its origin:
domain, source file, and line span.  Every query result carries this identifier.
When KGRAG output is passed to a synthesis model, the model receives facts with
addresses, not embeddings without provenance.

This makes synthesis auditable.  Every claim in a synthesized response can be
traced to a specific node in a specific graph, which can be traced to a specific
location in a primary source.  The model cannot introduce facts that are not in
the context, because the context is explicitly bounded by the retrieved graph
substructure.

% ============================================================
\section{Implications}
% ============================================================

\subsection{Training as Exception, Not Rule}

The haystack paradigm treats training as the central operation.  The forest
paradigm inverts this.  For strong-ontology domains, training is unnecessary.
For weak-ontology domains, training is localized to the reader: a small model
capable of query formulation and synthesis, not a massive compression of domain
knowledge.  Training is reduced from a universal knowledge compression problem to
a focused reading capability problem.

\subsection{Knowledge as Infrastructure}

In the haystack paradigm, knowledge is locked inside model weights.  Updating
knowledge requires retraining---expensive and opaque.  In the forest paradigm,
knowledge is in the graphs: explicit, structured, independently versionable.
Updating knowledge means adding nodes and edges.  No retraining; no parameter
updates; no risk of catastrophic forgetting.

Knowledge becomes an infrastructure asset rather than a training artifact---
inspectable, auditable, versioned, backed up, and incrementally updated like any
database.

\subsection{Correctness by Construction}

The most important property of the forest paradigm is one the haystack paradigm
cannot offer: correctness by construction at the knowledge layer.  Every edge in
a structurally derived knowledge graph is a theorem about its source.  There is
no extraction step where errors can be silently introduced by a probabilistic
model.

Synthesis can still produce reasoning errors.  Semantic search can return
less-relevant entry points.  Graph coverage is always incomplete relative to the
full domain.  But none of these errors are hallucinations---they are bounded,
auditable, and improvable through targeted corpus expansion and traversal
refinement.

The haystack hallucinates because it has no ground.  The forest does not
hallucinate because every tree has roots.

% ============================================================
\section{Conclusion}
% ============================================================

We have built haystacks because we believed that knowledge could be universally
compressed, and that a large enough compression would recover any fact on demand.
Those successes have run into the structural limits of the paradigm:
hallucination, ontological entanglement, opaque provenance, and the relentless
cost of retraining.

The alternative is not a better haystack.  For domains with strong ontologies,
the knowledge graph \emph{is} the knowledge---the corpus, compiled into a
traversable graph, contains answers as paths waiting to be walked.  What is
needed is not a more powerful model.  It is a larger corpus and a better reader.
For domains without strong ontologies, the direction is the same: impose
structure through natural-language tools, grow the graph incrementally, and let
provisional structure harden into formal ontology over time.

KGRAG implements this vision as a federated forest of knowledge trees.  Each
domain contributes its tree.  The federation layer makes the whole forest
queryable through a single interface.  The model receives structured,
provenance-tagged facts and synthesizes over them---a task for which small,
focused models are entirely adequate.

The forest is not a metaphor.  It is an architecture.  It is available today,
for the domains where formal ontologies already exist.  The remaining work is
construction: growing the corpus, expanding the domains, and building the readers
that can navigate the trees.

Every haystack we stop building is a forest we can start growing.

% ============================================================
\section*{Acknowledgments}
% ============================================================

The author thanks the open-source communities behind SQLite, LanceDB,
sentence-transformers, and the Python AST infrastructure that make deterministic
knowledge compilation practical.

% ============================================================
\bibliographystyle{plain}
\begin{thebibliography}{9}

\bibitem{kanehisa2023kegg}
M.~Kanehisa, M.~Furumichi, Y.~Sato, M.~Ishiguro-Watanabe, and M.~Tanabe.
\newblock {KEGG} for taxonomy-based analysis of pathways and genomes.
\newblock \emph{Nucleic Acids Research}, 51(D1):D587--D592, 2023.

\bibitem{shazeer2017moe}
N.~Shazeer, A.~Mirhoseini, K.~Maziarz, A.~Davis, Q.~Le, G.~Hinton, and
  J.~Dean.
\newblock Outrageously large neural networks: The sparsely-gated mixture-of-experts layer.
\newblock In \emph{International Conference on Learning Representations (ICLR)}, 2017.

\bibitem{lewis2020rag}
P.~Lewis, E.~Perez, A.~Piktus, F.~Petroni, V.~Karpukhin, N.~Goyal,
  H.~K{\"u}ttler, M.~Lewis, W.-T. Yih, T.~Rockt{\"a}schel, S.~Riedel, and
  D.~Kiela.
\newblock Retrieval-augmented generation for knowledge-intensive {NLP} tasks.
\newblock In \emph{Advances in Neural Information Processing Systems (NeurIPS)},
  2020.

\bibitem{edge2024graphrag}
D.~Edge, H.~Trinh, N.~Cheng, J.~Bradley, A.~Chao, A.~Mody, S.~Truitt, and
  J.~Larson.
\newblock From local to global: A {GraphRAG} approach to query-focused
  summarization.
\newblock \emph{arXiv preprint arXiv:2404.16130}, 2024.

\bibitem{bender2021parrots}
E.~M.~Bender, T.~Gebru, A.~McMillan-Major, and S.~Shmitchell.
\newblock On the dangers of stochastic parrots: Can language models be too big?
\newblock In \emph{ACM Conference on Fairness, Accountability, and Transparency
  (FAccT)}, pages 610--623, 2021.

\end{thebibliography}

\end{document}
